\chapter{Git と GitHub の基本}
\sheet{5-1〜5-6}

\section{Git がないと何が困るか}

レポートを書いていて、こんなファイルが並んだことはありませんか。

\begin{center}\ttfamily\small
\begin{tabular}{l}
レポート.docx\\
レポート\_修正.docx\\
レポート\_修正2.docx\\
レポート\_最終版.docx\\
レポート\_最終版\_v2\_本当に最終.docx
\end{tabular}
\end{center}

どれが最新なのか、何を変えたのか、どこで壊れたのかが分からなくなります。
これを人数分・ファイル数百個分でやるのが Web 開発です。
\term{Git}（ギット）は、ファイルを上書きしながら、\textbf{いつ・誰が・何を・なぜ変えたか}を全部記録してくれる道具です。

\section{Git の3つの場所}

Git を使うとき、ファイルは次の4つの場所を行き来します。最初はこの図だけ覚えれば十分です。

\begin{center}
\begin{tikzpicture}[font=\sffamily\small,
  area/.style={draw,rounded corners=3pt,minimum width=2.9cm,minimum height=2.3cm,align=center,fill=white},
  arr/.style={-{Stealth[length=2.5mm]},thick}]
  \node[area,draw=kincha] (w) {\textbf{作業場}\\{\scriptsize ワーキングツリー}\\[3pt]{\scriptsize いま編集している}\\{\scriptsize ファイル}};
  \node[area,draw=kon,right=1.1cm of w] (s) {\textbf{準備エリア}\\{\scriptsize ステージ}\\[3pt]{\scriptsize 次に記録する}\\{\scriptsize 予定のもの}};
  \node[area,draw=matcha,right=1.1cm of s] (r) {\textbf{記録}\\{\scriptsize ローカルリポジトリ}\\[3pt]{\scriptsize 自分の PC の}\\{\scriptsize 履歴}};
  \node[area,draw=gray,right=1.1cm of r] (g) {\textbf{GitHub}\\{\scriptsize リモートリポジトリ}\\[3pt]{\scriptsize クラウドの}\\{\scriptsize 履歴}};
  \draw[arr,kincha] (w) -- node[above,font=\ttfamily\scriptsize]{add} (s);
  \draw[arr,kon] (s) -- node[above,font=\ttfamily\scriptsize]{commit} (r);
  \draw[arr,matcha] (r) -- node[above,font=\ttfamily\scriptsize]{push} (g);
  \draw[arr,gray] (g.south) to[bend left=25] node[below,font=\ttfamily\scriptsize]{pull / clone} (w.south);
  \node[draw=gray!50,dashed,rounded corners,fit=(w)(s)(r),inner sep=6pt,label={[font=\scriptsize]above:あなたの PC（ローカル）}] {};
\end{tikzpicture}
\end{center}

\begin{description}[font=\sffamily\bfseries]
  \item[add] 「このファイルの変更を次の記録に入れる」と選ぶ
  \item[commit] 選んだ変更を、メッセージ付きで1つの記録（\term{コミット}）にする。ゲームのセーブに近い
  \item[push] 自分の PC にある記録を GitHub に送る
  \item[pull / clone] GitHub にある記録を自分の PC に持ってくる
\end{description}

\begin{why}[なぜ add と commit が分かれているのか]
1回の作業で「トップページの文言修正」と「FAQ の追加」を同時にやってしまったとします。
add で文言修正のファイルだけを選んで commit し、次に FAQ のファイルを選んで commit すれば、
\textbf{1つのコミットに1つの意味}を持たせられます。
あとで「FAQ の追加だけ取り消したい」となったときに、これが効いてきます。
\end{why}

\section{はじめてのリポジトリ}

\begin{step}{GitHub にリポジトリを作る}
\sheet{5-2}
\url{https://github.com/new} を開き、次のように入力して \btn{Create repository} を押します。
\screenshot[0.85]{github-new-repo.jpg}
\begin{itemize}
  \item \menu{Repository name}：\code{git-practice}
  \item \menu{Choose visibility}：\textbf{Public}
  \item \menu{Add README}：\textbf{Off のまま}（README は自分で作るため）
\end{itemize}
作ると、\code{…or push an existing repository from the command line} という説明が出た画面になります。
この画面は開いたままにしておきます。
\end{step}

\begin{step}{ローカルにリポジトリを作る}
\begin{terminal}[Git Bash]
$ cd ~/dev
$ mkdir git-practice
$ cd git-practice
$ git init
Initialized empty Git repository in C:/Users/ryusei/dev/git-practice/.git/
$ ls -a
.  ..  .git
\end{terminal}
\code{.git} という隠しフォルダが、Git の記録帳です。\textbf{このフォルダは手で触らないでください。}
\end{step}

\begin{step}{README を書く}
\begin{terminal}[Git Bash]
$ code README.md
\end{terminal}
VS Code が開くので、次の内容を書いて \key{Ctrl}+\key{S} で保存します。
\begin{filebox}{README.md}
\begin{Piton}[language=code]
# git-practice

Git を練習するためのリポジトリです。

## 自己紹介

- 名前：柴山りゅうせい
- 所属：理工学部 応用化学科 2年／理工体剣道部
- 好きな技：（ここに書く）
\end{Piton}
\end{filebox}
\code{.md} は \term{Markdown}（マークダウン）という書式です。
\code{\#} が見出し、\code{-} が箇条書きになります。GitHub ではきれいに整えて表示されます。
\end{step}

\begin{step}{status → add → commit}
\begin{terminal}[Git Bash]
$ git status
On branch main
No commits yet
Untracked files:
  (use "git add <file>..." to include in what will be committed)
        README.md

$ git add README.md
$ git status
Changes to be committed:
        new file:   README.md

$ git commit -m "docs: README を追加"
[main (root-commit) 3f2a1c9] docs: README を追加
 1 file changed, 9 insertions(+)
\end{terminal}
\textbf{\code{git status} は何度打ってもかまいません。}迷ったらまず \code{git status} です。
いまどの場所に何があるのかを教えてくれます。
\end{step}

\begin{step}{GitHub に push する}
さっき開いておいた GitHub の画面にある \textbf{SSH} のアドレス（\code{git@github.com:...}）を使います。
\begin{terminal}[Git Bash]
$ git remote add origin git@github.com:shibayama-ryusei/git-practice.git
$ git push -u origin main
...
To github.com:shibayama-ryusei/git-practice.git
 * [new branch]      main -> main
\end{terminal}
GitHub の画面を再読み込みして、README が表示されていれば成功です。
\code{origin} は「GitHub 側のリポジトリ」につけたあだ名です。
\code{-u} は初回だけ付ければよく、2回目からは \code{git push} だけで送れます。
\end{step}

\section{コミットメッセージの書き方}

この教材と本番のリポジトリでは、コミットメッセージを次の形で書きます（\term{Conventional Commits}）。

\begin{center}\small
\begin{tabularx}{\linewidth}{l l X}
\toprule
種類 & 使うとき & 例 \\
\midrule
\code{feat} & 新しい機能やページを足した & \code{feat: FAQ ページを追加} \\
\code{fix} & バグを直した & \code{fix: スマホで見出しがはみ出るのを修正} \\
\code{style} & 見た目だけを変えた & \code{style: ボタンの色を紺に変更} \\
\code{docs} & 説明文や README を変えた & \code{docs: README に起動方法を追記} \\
\code{refactor} & 動きは同じまま、書き方を整理した & \code{refactor: 部員表示を関数に切り出す} \\
\code{chore} & 設定などの雑務 & \code{chore: .gitignore を追加} \\
\bottomrule
\end{tabularx}
\end{center}

\begin{tip}[よいコミットの大きさ]
「1つのコミット＝説明が1行で済む変更」が目安です。
メッセージに「と」を書きたくなったら（「見出しを直したのと、FAQ を足した」）、コミットを分けるサインです。
こまめにコミットしておくと、壊したときに戻れる地点が増えます。
\end{tip}

\section{差分を見る}

\begin{step}{変更を加えて diff を見る}
README の「好きな技」を書き換えて保存してから、次を打ちます。
\begin{terminal}[Git Bash]
$ git diff
diff --git a/README.md b/README.md
@@ -7,4 +7,4 @@
 - 名前：柴山りゅうせい
 - 所属：理工学部 応用化学科 2年／理工体剣道部
-- 好きな技：（ここに書く）
+- 好きな技：出小手
\end{terminal}
\code{-} で始まる行が消えた行、\code{+} で始まる行が足された行です。
\textbf{\code{git diff} を読む力は、AI と開発するときに一番大事な力です。}
Claude Code が何を変えたのかは、最終的にはこれで確かめます。
\end{step}

\begin{step}{履歴を見る}
変更を add・commit・push してから、履歴を見ます。
\begin{terminal}[Git Bash]
$ git add README.md
$ git commit -m "docs: 好きな技を追記"
$ git push
$ git log --oneline
a81d0e4 (HEAD -> main, origin/main) docs: 好きな技を追記
3f2a1c9 docs: README を追加
\end{terminal}
\end{step}

\begin{tip}[VS Code でも差分は見られる]
VS Code の左端にある、枝分かれしたアイコン（ソース管理）を押すと、変更したファイルの一覧が出ます。
ファイル名を押すと、変更前と変更後が左右に並んで表示されます。
ターミナルの \code{git diff} と同じものなので、見やすい方を使ってかまいません。
\end{tip}

\section{取り消し方}

\begin{center}\small
\begin{tabularx}{\linewidth}{X l}
\toprule
状況 & コマンド \\
\midrule
保存したけど、まだ add していない変更をなかったことにしたい & \code{git restore ファイル名} \\
add したけど、やっぱり準備エリアから外したい & \code{git restore --staged ファイル名} \\
直前のコミットのメッセージを直したい（\textbf{push する前だけ}） & \code{git commit --amend -m "新しいメッセージ"} \\
\bottomrule
\end{tabularx}
\end{center}

\begin{caution}[push した後の履歴は書き換えない]
\code{--amend} や \code{git reset} で履歴を書き換えたあとに push しようとすると、拒否されます。
そこで \code{git push -f}（強制 push）を使うと、GitHub 側の履歴を上書きしてしまい、
チーム開発では他の人の作業を消す事故になります。
\textbf{強制 push は、三角に相談してからでないと使わない}と決めておいてください。
\end{caution}

\section{.gitignore：Git に記録させないもの}

パスワードや、自動で作られる大量のファイルは、Git に記録してはいけません。
記録させないファイルを書いておくのが \code{.gitignore} です。

\begin{filebox}{.gitignore（第2巻以降で使う例）}
\begin{Piton}[language=code]
# 自動で作られる大量のファイル
node_modules/
.next/

# 秘密の設定（API キーなど）
.env*
!.env.example

# OS が勝手に作るファイル
Thumbs.db
.DS_Store
\end{Piton}
\end{filebox}

\begin{caution}[秘密情報を push してしまったら]
API キーやパスワードを含むファイルを GitHub に push してしまったら、\textbf{すぐに三角に連絡してください}。
コミットを消しても、一度公開されたキーは誰かに見られた前提で扱う必要があり、キーそのものを作り直す必要があります。
隠さずにすぐ言ってくれれば、被害はほぼ防げます。
\end{caution}

\section{よくあるエラー}

\begin{center}\small
\begin{tabularx}{\linewidth}{>{\ttfamily\raggedright\arraybackslash}p{5.2cm} X}
\toprule
\normalfont エラーの一部 & 原因と対処 \\
\midrule
Permission denied (publickey) & SSH 鍵が GitHub に登録されていないか、HTTPS のアドレスと混ざっている。\code{ssh -T git@github.com} で確かめる \\
\addlinespace
Please tell me who you are & 初期設定（\code{user.name}、\code{user.email}）をしていない。第3章に戻る \\
\addlinespace
nothing to commit, working tree clean & 記録する変更がない。保存（\key{Ctrl}+\key{S}）し忘れていないか確認する \\
\addlinespace
! [rejected] ... (fetch first) & GitHub 側に、自分の PC にない記録がある。\code{git pull} してから push する \\
\addlinespace
not a git repository & Git の管理下にないフォルダにいる。\code{pwd} で場所を確認する \\
\addlinespace
\normalfont（\code{-m} を付け忘れて VS Code のタブが開いた） & メッセージ入力用のタブ。1行目にメッセージを書いて保存し、タブを閉じれば commit される \\
\bottomrule
\end{tabularx}
\end{center}

\begin{challenge}{わざと失敗して、元に戻す}
\sheet{5-6}
\code{git-practice} で次のことを順にやり、各段階で \code{git status} の表示がどう変わるかをメモしてください。
\begin{enumerate}
  \item README を適当に書き換えて保存する → \code{git restore} で元に戻す
  \item \code{keiko.md} という新しいファイルを作り、add する → \code{git restore --staged} で準備エリアから外す
  \item \code{keiko.md} を commit する。そのときメッセージをわざと間違える → \code{--amend} で直してから push する
\end{enumerate}
「どの場所にあるファイルを、どの場所に戻したのか」を、本章の3つの場所の図を使って定例で説明します。
\end{challenge}

\begin{checklist}
\begin{itemize}
  \checkitem 作業場・準備エリア・ローカルの記録・GitHub の4つの関係を図で説明できる
  \checkitem \code{status → add → commit → push} を、資料を見ずにできる
  \checkitem \code{git diff} の \code{+} と \code{-} を読める
  \checkitem Conventional Commits の形でメッセージを書ける
  \checkitem 強制 push をしてはいけない理由を言える
  \checkitem 秘密情報を push してしまったときの対応を言える
\end{itemize}
\end{checklist}
